% Value of the Data (mandatory bulleted list). The master supplies the
% \section*{Value of the Data} header; this fragment is only the itemize.
% All numbers MUST match paper/_shared/CANONICAL_FACTS.md and the CSVs.
% Descriptive only — no scientific conclusions (those belong to the parent
% JSR research article \cite{yu2026jsr}).
\begin{itemize}

  \item The database is, to the authors' knowledge, the first openly published,
        externally curated apogee benchmark for supersonic sounding rockets and
        high-power rocketry: the 25-flight headline corpus is drawn from a
        third-party RASAero~II altitude-comparison collection \cite{rogers2015},
        so flight selection is independent of any outcome of the simulator
        described in the parent article \cite{yu2026jsr}.

  \item Each flight is supplied with a paired commercial-baseline prediction
        (RASAero~II \cite{rogers2015}) alongside the measured apogee, enabling a
        head-to-head, apples-to-apples evaluation of any rocket-trajectory
        simulator against both ground truth and an established reference tool
        without re-deriving the baseline.

  \item Every reported statistic is reproducible: the published analysis scripts
        regenerate all signed errors, summary statistics, and bootstrap
        confidence intervals directly from the full-precision apogee columns, so
        reusers can audit and re-derive every figure rather than trusting
        pre-computed values.

  \item The corpus spans the subsonic, transonic, and supersonic regimes from
        Mach~0.54 to Mach~4.33, with an exploratory extension of roughly 20
        historical high-Mach sounding-rocket flights reaching toward Mach~7
        (e.g.\ Black Brant~V~VB at Mach~7.22), giving a continuous Mach ladder
        for regime-by-regime assessment.

  \item The data are reusable for validating or benchmarking any rocket
        aerodynamics or trajectory tool---including open-source codes such as
        OpenRocket \cite{niskanen2009}, RocketPy \cite{rocketpy2021}, and
        engineering-level prediction methods \cite{sooy2005}---against a common,
        version-locked set of flights.

  \item Each record carries vehicle, motor, body-diameter, launch-site altitude,
        and instrumentation-type metadata, together with the cited source
        document, so reusers can reconstruct and re-simulate the flights or trace
        every value back to its original published record.

\end{itemize}
